\begin{proposition}\label{mum}
The resulting fillings have smooth total spaces and reduced irreducible
central fibers. At $p_1,p_2$, the normalization is
$\PP_E(\cO_E\oplus L)$, with $E$ elliptic and $L\in\Pic^0(E)$; its two
boundary sections are identified by a translation. At $p_3$, the
normalization is $\dP_6$, with opposite boundary curves identified.
The respective Euler characteristics are $0,0,2$.

The section $e$ extends through the smooth locus of each filling.
Let $W$ be a central fiber and $F$ a nearby torus. The monodromy
$T:H^1(F,\ZZ)\to H^1(F,\ZZ)$ satisfies $(T-I)^2=0$, and its image
$(T-I)H^1(F,\ZZ)$ is a primitive sublattice, of rank one at
$p_1,p_2$ and rank two at $p_3$. Specialization induces isomorphisms
\begin{align}\label{sp}
 H^q(W,\ZZ)\xrightarrow{\ \simeq\ }H^q(F,\ZZ)^T,
 \qquad 0\leq q\leq4,
\end{align}
where in degree $q$ we use $T$ for the induced monodromy on
$H^q(F,\ZZ)$. Specialization on fundamental groups is surjective, with kernel
generated by the primitive vanishing cycles.
\end{proposition}
